% !TEX root = ../print.tex
% Draft \S2 --- Preliminaries and notation
%
% Source: draft/spatial-hemigroup-scale-space.md, \S2 (lines 12-138).
%
% This is the chapter where the trust boundary is set. Four of the draft's numbered
% clauses are classical Fourier analysis and are taken as cited interfaces: Bochner,
% Schoenberg, the symmetric \Levy--Khintchine representation with its unique triple, and
% \Levy's continuity theorem. Two further facts the draft uses without stating them --
% uniqueness of a Fourier--Stieltjes transform, and the continuity theorem in its
% measure form -- are recorded here as nodes so that the trust they carry is visible in
% the dependency graph rather than buried inside proofs; that is the causal article's
% convention and its reason.
%
% Three things the draft names as candidates for the ledger are NOT taken on trust, and
% the reroute is recorded at each node: the quadratic growth bound for continuous
% negative definite functions (\ref{lem:quadratic-growth}, proved from the truncation
% split in three lines, so Jacob Vol. I Lemma 3.6.22 is not needed); the nonvanishing of
% infinitely divisible transforms (Sato Lemma 7.5, superseded by the elementary
% \ref{lem:nonvanishing} of Chapter 4); and the null-array convergence theorem (Sato
% Thm. 8.7, superseded by the self-contained proof of \ref{thm:increments-levy}).

\chapter{Preliminaries and notation}
\label{ch:preliminaries}

Space runs over $\RR$ with coordinate $x$ and the transform variable is the wavenumber
$\omega$; scale is $t$, and a stage is indexed by an ordered pair of scales $s \le t$, with
a third letter $r \le s \le t$ when three are needed. The dilation ratio is $\lambda$.

We write $\Lone = L^1(\RR, dx)$, real, with positive cone $L^1_+$ and
$\int f := \int_\RR f(x)\,dx$; $M_+(\RR)$ for the finite positive Borel measures on $\RR$,
of total mass $\|\mu\| = \mu(\RR)$; translations $(\Tr{a} f)(x) = f(x-a)$; the reflection
$(Rf)(x) = f(-x)$, extended to measures as the pushforward under $x \mapsto -x$, a measure
being \emph{symmetric} when $R\mu = \mu$; dilations $(\Dil{\lambda} f)(x) = \lambda^{-1}
f(x/\lambda)$ on functions and the pushforward under $x \mapsto \lambda x$ on measures; and
convolution $(\mu * f)(x) = \int f(x-y)\,\mu(dy)$, for which $\|\mu * f\|_1 \le \|\mu\|\,\|f\|_1$,
$\Dil{\lambda}(\mu * f) = (\Dil{\lambda}\mu) * (\Dil{\lambda}f)$ and
$R(\mu * f) = (R\mu) * (Rf)$.

For $\mu \in M_+(\RR)$ the Fourier transform is
\begin{equation}\label{eq:transform}
  \hat\mu(\omega) = \int_\RR e^{-i\omega x}\,\mu(dx), \qquad \omega \in \RR,
\end{equation}
continuous, with $|\hat\mu| \le \|\mu\|$ and $\hat\mu(0) = \|\mu\|$; for symmetric $\mu$ it is
real and even, $\hat\mu(\omega) = \int \cos(\omega x)\,\mu(dx)$. Under dilation
$\widehat{\Dil{\lambda}\mu}(\omega) = \hat\mu(\lambda\omega)$, and
$\widehat{\mu * \nu} = \hat\mu\,\hat\nu$. Unlike the Laplace transform of a measure carried by
a half-line, $\hat\mu$ need not be positive; that the axioms exclude this is
\ref{lem:nonvanishing}.

% draft: Definition 2.1
% CHANGED (review 2026-09-09): the inequality now says "\ge 0 (in particular real)". The
% displayed sum is complex, and the standard reading -- the only one under which
% prop:fourier-toolbox(1) is true -- is that it is a NONNEGATIVE REAL, so that Hermitian
% symmetry of varphi is part of the assertion rather than a separate hypothesis. The Lean
% (SpatialLine.IsPositiveDefinite) reads it that way, through 0 <= z under ComplexOrder;
% answering Q4, the prose should commit to it too.
\begin{definition}[positive definite]\label{def:positive-definite}
\lean{SpatialLine.IsPositiveDefinite}\leanok
\notes{positive-definite-functions}
A function $\varphi : \RR \to \CC$ is \emph{positive definite} if
\[
  \sum_{j,k} c_j\,\overline{c_k}\,\varphi(\omega_j - \omega_k) \ge 0
  \quad\text{(in particular real)}
\]
for every finite family $(\omega_j) \subset \RR$ and every finite family $(c_j) \subset \CC$.
\statusT\quad\emph{(The parenthesis is part of the definition: the double sum of a complex
function is asserted to be a nonnegative real, which is what makes $\varphi$ Hermitian ---
$\varphi(-\omega) = \overline{\varphi(\omega)}$ --- and it is the reading under which
\ref{prop:fourier-toolbox}(1) holds.)}
\end{definition}

% draft: Definition 2.2
% twin: hcs def:bernstein-function (Hemigroup.levyExponent) --- the causal class is defined
% by complete monotonicity of the derivative, this one by positive definiteness of the
% exponential; the two definitions play the same structural role.
\begin{definition}[symmetric continuous negative definite function]\label{def:symmetric-negdef}
\uses{def:positive-definite}
\lean{SpatialLine.IsSymNegDef}\leanok
\notes{positive-definite-functions}
A function $\psi : \RR \to [0,\infty)$ is a \emph{symmetric continuous negative definite
function} if $\psi$ is continuous and even, $\psi(0) = 0$, and $e^{-\tau\psi}$ is positive
definite for every $\tau > 0$. We write $\psi \in \NDs$.
\statusT\quad\emph{(The exponential form is taken as primitive, so that the class is defined
without any derivative and Schoenberg's theorem is needed only to import results stated for
the kernel form. This is the choice the causal article made for its Bernstein class in the
opposite direction, and it is why \ref{prop:fourier-toolbox}(2) is an interface rather than a
definition.)}
\end{definition}

% draft: Proposition 2.3
% twin: hcs prop:bernstein-toolbox (ledger A1-A4 there); the four clauses correspond one to
% one, with Bochner in the place of Bernstein--Widder and Schoenberg in the place of
% "e^{-tau g} completely monotone for every tau".
% NOTE (proof 2026-09-09): the CONVERSE half of clause (3) is now admitted in the Lean as the
% axiom SpatialLine.fourier_toolbox_levy_converse, on blueprint/trust-boundary.txt, with A3's
% **Lean:** line naming it. Its only consumer is lem:profile-integrability's "consequently"
% clause. The other clauses are admitted nowhere; the statement is unchanged.
\begin{proposition}[the Fourier toolbox]\label{prop:fourier-toolbox}
\uses{def:positive-definite,def:symmetric-negdef}
\lean{Skeleton.fourier_toolbox_bochner, Skeleton.fourier_toolbox_bochner_symm,
Skeleton.fourier_toolbox_schoenberg, Skeleton.fourier_toolbox_levy_exists,
Skeleton.fourier_toolbox_levy_converse, Skeleton.fourier_toolbox_levy_unique,
Skeleton.fourier_toolbox_cone, Skeleton.fourier_toolbox_kernel_closed,
Skeleton.fourier_toolbox_closure}\notready
\notes{positive-definite-functions}
\begin{enumerate}
\item \emph{(Bochner.)} $\varphi$ is the Fourier transform of a probability measure on $\RR$
if and only if $\varphi$ is continuous, positive definite and $\varphi(0) = 1$; it is the
transform of a \emph{symmetric} probability measure if and only if moreover $\varphi$ is real,
equivalently even.
\item \emph{(Schoenberg.)} For continuous even $\psi \ge 0$ with $\psi(0) = 0$: $\psi \in \NDs$
if and only if $\psi$ is negative definite in the kernel sense,
\[
  \sum_{j,k}\bigl(\psi(\omega_j) + \psi(\omega_k) - \psi(\omega_j - \omega_k)\bigr)\,
  c_j\,\overline{c_k} \ \ge\ 0
\]
for every finite family; equivalently, by clause (1), if and only if $e^{-\tau\psi}$ is the
transform of a symmetric probability measure for every $\tau > 0$.
\item \emph{(\Levy--Khintchine, symmetric form.)} Every $\psi \in \NDs$ has the representation
\begin{equation}\label{eq:levy-khintchine}
  \psi(\omega) \;=\; a\,\omega^2 + \int_{(0,\infty)} \bigl(1 - \cos\omega x\bigr)\,\nu(dx),
  \qquad a \ge 0, \quad \int_{(0,\infty)} (1 \wedge x^2)\,\nu(dx) < \infty,
\end{equation}
the pair $(a,\nu)$ being unique; and conversely every such pair defines a member of $\NDs$.
\item \emph{(Closure.)} $\NDs$ is a convex cone, and in the kernel form of clause (2) the
class is closed under pointwise limits with no proviso. If $\psi_n \in \NDs$ and
$\psi_n \to \psi$ pointwise with $\psi$ continuous at $0$, then $\psi \in \NDs$.
\end{enumerate}
\statusA\quad\ledger{A1}\quad\ledger{A2}\quad\ledger{A3}\quad\ledger{A4}\quad\emph{(The
classical Fourier analysis this article stands on, taken as four cited interfaces and not
reproven. \textbf{Assignment.} \ledger{A1} carries clause (1), both the general and the
symmetric leg; \ledger{A2} carries clause (2); \ledger{A3} carries clause (3) --- existence,
converse, and uniqueness of the pair; \ledger{A4} carries clause (4), the cone and both
closure statements. What none of them carries is the \emph{folding} convention by which a
symmetric \Levy{} measure on $\RR \setminus \{0\}$ is written as a measure on $(0,\infty)$: the
sources state clause (3) on $\RR$ with a centering term, and the reduction to
\ref{eq:levy-khintchine} --- the centering term vanishes by symmetry, and the two halves of
the measure are added under $x \mapsto |x|$ --- is ours and elementary, written out
immediately below. The classes $\NDs$ and $\LEs$ (\ref{def:levy-exponents}) are the same class,
named by the definition and by the representation respectively, and the clause that identifies
them is (3) of this proposition; this article's proofs manipulate the representation, and the
machine-checked development defines $\LEs$ alone.)}
\end{proposition}

% ADDED (article mirror 2026-09-14, R157): written for the paper at the author's request
% (2026-09-13), condensed from the hub's wiki/positive-definite-functions.md, and mirrored here
% because the blueprint is the text of record for what the paper prints. The Feller anchors are
% the ones the librarian confirmed (Lemma 1 p. 621 with the even window and the Gaussian factor
% (2.4); the Toeplitz form (2.6)-(2.7) p. 622; Lemma 3 p. 622).
\begin{remark}[what the two classes are, and why the transform sees them]
\label{rem:pd-nd-intuition}
The prototype of a positive definite function is a single \emph{character},
$\omega \mapsto e^{-i\omega x}$ for a fixed $x$: the double sum of \ref{def:positive-definite}
collapses to $|\sum_j c_j e^{-i\omega_j x}|^2 \ge 0$. The condition is linear in $\varphi$ and
survives pointwise limits, so every positive mixture of characters, that is every Fourier
transform of a positive measure, is positive definite for free. Bochner's theorem is the
converse: a continuous positive definite function is such a mixture, and the mixing measure is
unique. In the language of convex cones, the characters are the extreme rays of the cone of
positive definite functions, and every continuous member is the barycentre of the rays, with
the measure as the weights. In the symmetric case the rays are $\cos\omega x$, $x \ge 0$.

The negative definite functions are the same picture one level up. Ask which $\psi$ make
$e^{-\tau\psi}$ a transform for every $\tau > 0$, which is what the cascade of
Chapter~\ref{ch:cascade} will ask. The answer is again a cone described twice: by the checkable
kernel-form inequality of clause (2), and by its extreme rays, which is the \Levy--Khintchine
representation \ref{eq:levy-khintchine}. The rays are $\omega^2$ and $1 - \cos\omega x$ for
$x > 0$, and there is no constant ray because $\psi(0) = 0$. Each ray has a probabilistic
identity, given in \ref{rem:displacement-dictionary} below.

The mechanism behind Bochner's theorem is the one this article reuses. In Feller's proof the
transform is inverted against an even window and the general case is reached through the
convergence factor $e^{-\varepsilon\omega^2}$ (\sscite{feller2009introduction}, Vol.~2,
\S~XIX.2, Lemma~1, p.~621, and (2.4)); positive definiteness is what makes the windowed inverse
transforms nonnegative, and the finite-family inequality is the discrete form of that, the
Toeplitz form (2.6)--(2.7) on p.~622. The same device, a transform read against a Gaussian
window, is the test function of Chapter~\ref{ch:representation} and the smoothed transmittance
of Chapter~\ref{ch:cascade}. Two facts used throughout are also Feller's (Lemma~3, p.~622):
$|\varphi| \le \varphi(0)$ and $\varphi(-\omega) = \overline{\varphi(\omega)}$.

The contrast with the half-line is one of kind. Complete monotonicity, the transform-free
description of ``comes from positive mass'' for the Laplace transform, is a pointwise test:
countably many derivative signs at each point. Positive definiteness is a global quadratic-form
condition, and there is no local test for it. That is the first sign that the half-line theory
is the more rigid of the two.
\end{remark}

\emph{The folding convention.} A symmetric \Levy{} measure $\nu_2$ on $\RR \setminus \{0\}$
enters \ref{eq:levy-khintchine} as its image $\nu$ under $x \mapsto |x|$. So a two-sided
\Levy{} density written $h(|x|)/|x|$ --- the form the probabilistic sources use --- has folded
profile $k = 2h$ in \ref{eq:sd-profile} below: the factor is the two directions of a
displacement of a given size, and it is the one place where a constant is easily lost in
carrying a statement from a source into this article. The centering term of the general
representation vanishes by symmetry, read on the compensated integrand: the odd part of
$e^{-i\omega x} - 1 + i\omega x\,\mathbf 1_{|x| < 1}$ is
$-i(\sin\omega x - \omega x\,\mathbf 1_{|x|<1})$, which is $O(x^3)$ at the origin and bounded,
hence $\nu_2$-integrable, and odd, so its integral against a symmetric $\nu_2$ is $0$. The
uncompensated $\int \sin(\omega x)\,\nu_2(dx)$ need not converge.

\emph{Which vocabulary does the work.} The later chapters name the class by its representation.

% ADDED (article mirror 2026-09-14, R157): a node for the name the chapters from 5 on use. The
% blueprint had it only as a sentence of prop:fourier-toolbox's annotation, and the paper gave it
% a definition of its own at the author's request (2026-09-13); the blueprint is the text of
% record, so the node belongs here too. The identification LEs = NDs is prop:fourier-toolbox(3),
% ledger A3, and is stated in the annotation rather than in the definition, so that the \leanok
% covers exactly what SpatialLine.IsSymLevyExponent defines.
\begin{definition}[symmetric \Levy{} exponents]\label{def:levy-exponents}
\uses{prop:fourier-toolbox}
\lean{SpatialLine.IsSymLevyExponent}\leanok
$\LEs$ is the set of functions $\psi : \RR \to [0,\infty)$ of the form
\ref{eq:levy-khintchine}, for some $a \ge 0$ and some measure $\nu$ on $(0,\infty)$ with
$\int (1 \wedge x^2)\,\nu(dx) < \infty$.
\statusT\quad\emph{(The same class as $\NDs$ of \ref{def:symmetric-negdef}, named by its
representation instead of by its definition: that $\LEs = \NDs$ is \ref{prop:fourier-toolbox}(3)
and is cited there, on ledger A3. The article's proofs manipulate the representation, and the
machine-checked development defines $\LEs$ alone, so this is the name the chapters from
Chapter~\ref{ch:cascade} on use; \ref{def:symmetric-negdef} enters only where the class $\NDs$
is the subject, and in one leg of \ref{lem:selfdecomposable-exponents}.)}
\end{definition}

% additive: not in the draft, which uses the fact without stating it. Used by
% lem:convolution-representation, lem:additivity, thm:main-characterization.
% twin: hcs prop:laplace-uniqueness (ledger A6 there)
% CHANGED (proof 2026-09-09): DEMOTED from [A] (ledger A5) to [T]. The statement is not cited
% but proved: Mathlib's MeasureTheory.Measure.ext_of_charFun is it verbatim for finite measures
% on a complete second-countable inner product space, and SpatialLine.fourier_uniqueness is one
% line from it. Ledger A5 is retired -- kept for the record, grounding no node -- and no name
% reaches trust-boundary.txt. The author's decision of 2026-09-09, answering the statement
% review's Q2; the statement text itself is unchanged.
\begin{proposition}[uniqueness of the Fourier--Stieltjes transform]\label{prop:fourier-uniqueness}
\lean{SpatialLine.fourier_uniqueness}\leanok
A finite Borel measure on $\RR$ is determined by its Fourier transform: if
$\hat\mu(\omega) = \hat\rho(\omega)$ for every $\omega \in \RR$, then $\mu = \rho$.
\statusT\quad\emph{(Classical, and machine-checked rather than cited. The Lean declaration is
one line from Mathlib's \texttt{MeasureTheory.Measure.ext\_of\_charFun}, which is this statement
for finite measures on a complete second-countable inner product space, of which $\RR$ is one;
the node carried the grade \textbf{[A]} on ledger A5 until 2026-09-09, and A5 is now retired and
grounds no node. The statement is about \emph{finite} measures on $\RR$ and claims nothing for
infinite ones; the corresponding statement on $(0,\infty)$ is
\ref{prop:laplace-uniqueness-locally-finite}.)}
\end{proposition}

\begin{proof}\leanok
The functions $e_\omega(x) = e^{-i\omega x}$, $\omega \in \RR$, span a subalgebra $\mathcal{A}$
of the bounded continuous functions on $\RR$: it contains the constants, since $e_0 = 1$; it is
closed under conjugation, since $\overline{e_\omega} = e_{-\omega}$; and it separates points,
since $e_\omega(x) = e_\omega(y)$ for every $\omega$ forces $x = y$. The hypothesis says that
$\mu$ and $\rho$ give equal integrals to every $e_\omega$, hence, by linearity, to every member
of $\mathcal{A}$. Two finite Borel measures on a complete separable metric space that give equal
integrals to every member of such an algebra are equal: on each compact set the
Stone--Weierstrass theorem makes $\mathcal{A}$ uniformly dense in the continuous functions, and a
finite Borel measure on a Polish space is inner regular, so the equality extends first to the
bounded continuous functions and then, by the portmanteau theorem, to the Borel sets. This is
the argument Mathlib's \texttt{MeasureTheory.Measure.ext\_of\_charFun} runs, in the stated
generality; every step of it is machine-checked, and the article's declaration is its
specialisation to $\RR$.
\end{proof}

% additive: not in the draft, which uses the fact inside the proof of Theorem 7.3'. Used by
% thm:main-characterization, prop:matern-exponent, prop:thorin-machine.
% twin: hcs prop:laplace-continuity (ledger A5 there)
% CHANGED (proof 2026-09-09): DEMOTED from [A] (ledger A6) to [T], and SPLIT. The node now
% states only the clause the article consumes -- the limit is known in advance to be the
% transform of a probability measure -- which is proved, from Mathlib's
% MeasureTheory.ProbabilityMeasure.tendsto_of_tendsto_charFun. The general clause, in which the
% limit function is only continuous at the origin, is used by NO node: all four consumers
% (prop:two-members(1), thm:main-characterization, prop:matern-density, prop:thorin-machine)
% apply the theorem with the limit already in hand, which is what reading the \uses edges shows.
% It is therefore moved out of the statement into rem:levy-continuity-general below, and ledger
% A6 is retired -- kept for the record, grounding no node. No \uses edge changes: the label is
% the same and still names what every consumer needs. Author's decision of 2026-09-09,
% answering the statement review's Q2.
\begin{proposition}[\Levy's continuity theorem]\label{prop:levy-continuity}
\uses{prop:fourier-uniqueness}
\lean{SpatialLine.levy_continuity}\leanok
Let $\mu_n, \mu$ be probability measures on $\RR$. If $\hat\mu_n(\omega) \to \hat\mu(\omega)$
for every $\omega \in \RR$, then $\mu_n \to \mu$ weakly.
\statusT\quad\emph{(Classical, and machine-checked rather than cited; the node carried the grade
\textbf{[A]} on ledger A6 until 2026-09-09, and A6 is now retired and grounds no node. Weak
convergence is read, in the Lean statement, as convergence of $\int g\,d\mu_n$ for every bounded
continuous $g : \RR \to \RR$, which is the form the $\varepsilon/3$ argument in the proof of
\ref{thm:main-characterization} consumes; the Lean declaration is Mathlib's
\texttt{MeasureTheory.ProbabilityMeasure.tendsto\_of\_tendsto\_charFun} read through that
characterization of the weak topology. What is \emph{not} part of this statement is the passage
from weak convergence to convergence of the convolution operators in $\Lone$, which is the
$\varepsilon/3$ argument itself and is proved where it is used. The general form of the theorem,
in which the limit function is only assumed continuous at the origin, is
\ref{rem:levy-continuity-general}; no node of this article uses it.)}
\end{proposition}

\begin{proof}\leanok
Since $\hat\mu$ is continuous at $0$ with $\hat\mu(0) = 1$, the sequence is tight. Indeed, for a
probability measure $\rho$ on $\RR$ and $r > 0$, Fubini's theorem gives
\[
  \frac1r\int_{-r}^{r}\bigl(1 - \hat\rho(\omega)\bigr)\,d\omega
  = 2\int_\RR\Bigl(1 - \frac{\sin rx}{rx}\Bigr)\,\rho(dx)
  \ \ge\ \rho\bigl(\{|x| \ge 2/r\}\bigr),
\]
because the integrand is nonnegative and is at least $\tfrac12$ where $|rx| \ge 2$. Applying
this to $\mu_n$ and letting $n \to \infty$ --- the integrands are bounded by $2$, so bounded
convergence applies --- bounds the upper limit of $\mu_n(\{|x| \ge 2/r\})$ by
$\frac1r\int_{-r}^r(1 - \hat\mu)$, which tends to $0$ as $r \downarrow 0$ by continuity of
$\hat\mu$ at $0$. So $\{\mu_n\}$ is tight. By Prokhorov's theorem every subsequence has a
further subsequence converging weakly to some probability measure $\rho$; along it
$\hat\rho = \lim \hat\mu_n = \hat\mu$, since $x \mapsto e^{-i\omega x}$ is bounded and
continuous, so $\rho = \mu$ by \ref{prop:fourier-uniqueness}. Every subsequence therefore has a
further subsequence converging weakly to $\mu$, which is convergence of the whole sequence.
The statement is machine-checked: Mathlib's
\texttt{MeasureTheory.ProbabilityMeasure.tendsto\_of\_tendsto\_charFun} proves it with the same
tightness estimate (\texttt{isTightMeasureSet\_of\_tendsto\_charFun}) and, in place of the
subsequence extraction, a direct density argument in the algebra of characters.
\end{proof}

\begin{remark}[the general form of the continuity theorem]\label{rem:levy-continuity-general}
The classical statement is more general: if $\hat\mu_n \to \varphi$ pointwise with $\varphi$
continuous at $0$, then $\varphi$ is the transform of a probability measure $\mu$ and
$\mu_n \to \mu$ weakly. The tightness estimate in the proof above uses only continuity of the
limit function at the origin, so it applies unchanged. What the general form needs in addition
is that the limit along a Prokhorov subsequence has transform $\varphi$, which identifies
$\varphi$ as a transform. The continuity proviso is what the line costs against the half-line,
where pointwise limits of Laplace exponents on $(0,\infty)$ suffice --- and it is why
\ref{prop:fourier-toolbox}(4) carries it too. No statement of this article uses the general
form: every use of \ref{prop:levy-continuity} has the limit law already in hand.
\end{remark}

% additive: not in the draft. Used by thm:scale-monotone-noncreation and
% prop:thorin-subclass, where a \Levy{} tail has to be recovered from its Laplace transform.
% twin: hcs prop:laplace-uniqueness-sigma-finite
% (Hemigroup.laplaceL_injective_of_ne_top) --- proved there, so [T] here rather than [A].
\begin{proposition}[uniqueness of the Laplace transform of a locally finite measure]
\label{prop:laplace-uniqueness-locally-finite}
\lean{SpatialLine.laplace_uniqueness_locally_finite}\leanok
Let $m, m'$ be Borel measures on $(0,\infty)$, not necessarily finite, and suppose that
$\int e^{-\tau u}\,m(du) = \int e^{-\tau u}\,m'(du) < \infty$ for every $\tau$ in some
interval $(\tau_0,\infty)$. Then $m = m'$.
\statusT\quad\emph{(Elementary once finiteness of the transform is used to reduce to finite
measures. Held \textbf{[T]} rather than cited because its causal twin is proved in the causal
development and that argument transfers verbatim: the statement is about a measure on a
half-line and says nothing about $\RR$. Machine-checked 2026-09-09, on Lean core alone, as that
port; the Lean file mentions neither symmetry nor the Fourier transform, and the two
developments' copies now differ in one predicate, which is what makes it the shared core's first
candidate. The proof below uses finiteness of the transform at \emph{one} point of the ray, which
is weaker than the hypothesis stated and is the sharpening the causal side found.)}
\end{proposition}

\begin{proof}\leanok
Fix $\tau_1 > \tau_0$. Both transforms being finite at $\tau_1$, the measures
$e^{-\tau_1 u}m(du)$ and $e^{-\tau_1 u}m'(du)$ are finite, and they have equal Laplace
transforms on $(0,\infty)$ after the shift $\tau \mapsto \tau + \tau_1$; so it is enough to
treat finite measures. For finite $m$, the map $u \mapsto e^{-u}$ carries $(0,\infty)$
homeomorphically onto $(0,1)$, and $\int e^{-n u}\,m(du)$, $n \in \NN$, is the $n$-th moment
of the image measure on $[0,1]$. A finite measure on a bounded interval is determined by its
moments, the polynomials being dense in the continuous functions there. So the image measures
agree, and so do $m$ and $m'$.
\end{proof}

% additive: not in the draft, which cites a growth bound for continuous negative definite
% functions at two places (its Corollary 7.4' and its Theorem 12.2). REROUTED: the bound is
% three lines from the truncation split, so no ledger entry is opened for it.
\begin{lemma}[continuous negative definite functions grow at most quadratically]
\label{lem:quadratic-growth}
\uses{prop:fourier-toolbox}
\lean{SpatialLine.SymLevyPair.quadratic_growth,
SpatialLine.SymLevyPair.quadratic_growth_isBigO}\leanok
Let $\psi \in \NDs$ have the pair $(a,\nu)$ of \ref{eq:levy-khintchine}. Then
\[
  \psi(\omega) \ \le\ C\,(1 + \omega^2), \qquad \omega \in \RR,
  \qquad C := a + \tfrac12\int_{(0,1]} x^2\,\nu(dx) + 2\,\nu((1,\infty)) \ < \ \infty .
\]
In particular $\psi(\omega) = O(\omega^2)$ as $|\omega| \to \infty$.
\statusT\quad\emph{(Elementary, and it replaces a cited growth bound at both places the draft
used one. Both constants are finite exactly because $\int (1 \wedge x^2)\,\nu < \infty$.
Machine-checked 2026-09-09, on Lean core alone: the Lean statement is in $[0,\infty]$, so the
bound \emph{implies} finiteness of the exponent rather than presupposing it, which is what lets
the representation structure carry no finiteness field. The corollary
$\psi(\omega) < \infty$ is the helper \texttt{SymLevyPair.exponentL\_ne\_top}, on which every
statement of this article that reads an exponent as a real number rests.)}
\end{lemma}

\begin{proof}\leanok
Split the integral in \ref{eq:levy-khintchine} at $x = 1$. For $x \le 1$ use
$1 - \cos\omega x \le \tfrac12\omega^2x^2$; for $x > 1$ use $1 - \cos\omega x \le 2$. This
gives
$\psi(\omega) \le a\omega^2 + \tfrac12\omega^2\int_{(0,1]}x^2\,\nu(dx) + 2\,\nu((1,\infty))
\le C\,(1+\omega^2)$. Finiteness of the two integrals is
$\int (1 \wedge x^2)\,\nu(dx) < \infty$.
\end{proof}

The next statement is the draft's Remark 2.4, promoted to a node because three later proofs
turn on it. It is the exact point at which the line parts from the half-line: the causal
theory has a vanishing lemma saying that a member of its function class vanishing at one
interior point vanishes identically, and the Fourier-side analogue is false.

% draft: Remark 2.4 (promoted to a lemma: used by cor:monotonicity, lem:no-lattice,
% prop:strict-positivity and prop:bridge-strictness)
% twin: hcs lem:vanishing (Hemigroup.levyExponent_eq_zero_of_eq_zero) --- the twin is the
% statement this one contradicts: on the half-line the zero set of an exponent is trivial,
% on the line it is a closed subgroup and may be a lattice.
\begin{lemma}[the zero set of a symmetric exponent is a closed subgroup]\label{lem:lattice-zero}
\uses{prop:fourier-uniqueness,prop:fourier-toolbox}
\lean{SpatialLine.lattice_zero_isClosedSubgroup, SpatialLine.lattice_zero_trichotomy,
SpatialLine.lattice_zero_eq_univ_iff, SpatialLine.lattice_zero_mem_iff,
SpatialLine.lattice_zero_exponent, SpatialLine.lattice_zero_example_mem,
SpatialLine.lattice_zero_example_lattice}\leanok
\notes{positive-definite-functions}
Let $\mu$ be a probability measure on $\RR$ and put
$N(\mu) := \{\omega \in \RR : \hat\mu(\omega) = 1\}$. Then $N(\mu)$ is a closed subgroup of
$\RR$, hence $\{0\}$, a lattice $c\ZZ$ with $c > 0$, or $\RR$; the last case holds if and only
if $\mu = \delta_0$. For $\omega_0 \ne 0$,
\[
  \omega_0 \in N(\mu)
  \quad\Longleftrightarrow\quad
  \cos(\omega_0 x) = 1 \ \text{ for } \mu\text{-almost every } x
  \quad\Longleftrightarrow\quad
  \mu \text{ is carried by the lattice } \tfrac{2\pi}{|\omega_0|}\ZZ .
\]
If $\mu$ is symmetric with $\hat\mu = e^{-\psi}$ for some $\psi \in \NDs$, then
$N(\mu) = \{\psi = 0\}$: a symmetric exponent vanishes at a nonzero frequency exactly when its
law is a lattice law. Lattice laws can be infinitely divisible --- the symmetric compound
Poisson law on $\ZZ$, of exponent $c\,(1 - \cos\omega)$, is one --- so membership of $\NDs$
does not exclude them.
\statusT\quad\emph{(Elementary. The last sentence is what makes the node worth stating: it
records that no analogue of the causal vanishing lemma is available, and hence that the
exclusion of lattices must come from elsewhere. It comes from covariance,
\ref{lem:no-lattice}. Machine-checked 2026-09-09, all seven clauses, on Lean core alone. Two
readings the formal statement had to make and the prose need not: the almost-everywhere
quantifier in the displayed equivalences is over $\mu$ and not over Lebesgue measure, and
``carried by the lattice'' is $\mu$ of the complement being zero. The example's law is
\emph{quantified over} rather than constructed --- any probability measure with the stated
transform is carried by $\ZZ$ --- which is legitimate because such a measure is unique by
\ref{prop:fourier-uniqueness}.)}
\end{lemma}

\begin{proof}\leanok
Since $\Rea\bigl(1 - e^{-i\omega x}\bigr) = 1 - \cos\omega x \ge 0$, the identity
$\hat\mu(\omega) = 1$ forces $\int (1 - \cos\omega x)\,\mu(dx) = 0$, hence
$\cos(\omega x) = 1$ for $\mu$-almost every $x$, hence $e^{-i\omega x} = 1$ almost everywhere
and $\hat\mu(\omega) = 1$; the three conditions are therefore equivalent. So $N(\mu)$ is
closed under sums and negation, and closed as a set because $\hat\mu$ is continuous: a closed
subgroup of $\RR$, hence $\{0\}$, some $c\ZZ$ with $c > 0$, or $\RR$. If $N(\mu) = \RR$ then
$\hat\mu \equiv 1 = \hat\delta_0$, so $\mu = \delta_0$ by \ref{prop:fourier-uniqueness}; the
converse is immediate. For $\omega_0 \ne 0$, $\cos(\omega_0 x) = 1$ says
$x \in \tfrac{2\pi}{|\omega_0|}\ZZ$, so the almost-everywhere clause says exactly that $\mu$
is carried by that lattice. When $\hat\mu = e^{-\psi}$ with $\psi$ real, $\hat\mu(\omega) = 1$
if and only if $\psi(\omega) = 0$. For the example, $c\,(1-\cos\omega)$ is of the form
\ref{eq:levy-khintchine} with $a = 0$ and $\nu = c\,\delta_1$, hence lies in $\NDs$ by
\ref{prop:fourier-toolbox}(3); its law is carried by $\ZZ$, and it is infinitely divisible
because $c\,(1-\cos\omega)/n$ is again of that form.
\end{proof}

% draft: Definition 2.5
\begin{definition}[self-decomposable; class $L$]\label{def:self-decomposable}
\lean{SpatialLine.IsSelfDecomposable}\leanok
\notes{self-decomposability}
A probability measure $\mu$ on $\RR$ is \emph{self-decomposable} if for every $b > 1$ there is
a probability measure $\rho_b$ on $\RR$ with
$\hat\mu(\omega) = \hat\mu(\omega/b)\,\hat\rho_b(\omega)$ for every $\omega$; equivalently, if
$X \seq cX + R_c$ with $R_c$ independent of $X$, for every $c \in (0,1)$, where $X$ has law
$\mu$.
\statusT
\end{definition}

% draft: Proposition 2.6
% twin: hcs lem:selfdecomposable-exponents (ledger A18 there). The relation is deliberately
% NOT a Lean reuse: the causal article consumes its classical statement, and this article
% does not consume this one -- lem:selfdecomposable-exponents below proves the profile form
% it needs from the dilation identity instead.
% CHANGED (2026-09-11, draft sync, R1): the profile is now stated NONNEGATIVE where the class is
% introduced -- "k : (0,infty) -> [0,infty) nonincreasing" in the sentence and "k >= 0
% nonincreasing" at eq:sd-profile. Nonnegativity was implied by reading k(x)x^{-1}dx as a \Levy{}
% measure, and it was stated at lem:profile-integrability, but not at the display every later
% node refers to as "the form eq:sd-profile"; the Lean carries it as a field (SDProfile.k_nonneg).
% Nothing is narrowed: a nonincreasing k with a negative value is negative from that point on and
% gives no \Levy{} measure. Mirrored in the draft's (2.2) and Proposition 2.6 in the same sync.
\begin{proposition}[symmetric self-decomposable exponents]\label{prop:sd-exponents}
\uses{def:self-decomposable,prop:fourier-toolbox}
\lean{Skeleton.sd_exponents}\notready
\notes{self-decomposability}
A symmetric infinitely divisible law with pair $(a,\nu)$ is self-decomposable if and only if
its \Levy{} measure has the form $\nu(dx) = k(x)\,x^{-1}dx$ on $(0,\infty)$ with
$k : (0,\infty) \to [0,\infty)$ nonincreasing, so that its exponent is
\begin{equation}\label{eq:sd-profile}
  \psi(\omega) \;=\; a\,\omega^2
  + \int_0^\infty \bigl(1 - \cos\omega x\bigr)\,\frac{k(x)}{x}\,dx,
  \qquad a \ge 0, \quad k \ge 0 \ \text{nonincreasing}.
\end{equation}
Self-decomposability imposes no restriction on the Gaussian coefficient $a$.
\statusA\quad\ledger{A7}\quad\emph{(The classical statement, cited for orientation.
\textbf{Assignment.} \ledger{A7} carries the whole equivalence. What it does not carry is
anything this article's development consumes: the profile form is obtained independently in
\ref{lem:selfdecomposable-exponents}, from the dilation identity and a convex-tail argument,
and every later node cites that node and not this one. The node is kept because the statement
is part of the article's text of record and is where a reader meets the class; that it lies on
no proof path is deliberate, and is the difference from the causal article, which consumes
its twin.)}
\end{proposition}

% additive: the integrability criterion the draft records after its Proposition 2.6 and then
% cites repeatedly as "the criterion". Used by lem:selfdecomposable-exponents, lem:cin-rays,
% prop:choquet-cone, prop:matern-exponent.
% twin: hcs prop:admissibility-criterion + lem:criterion-converse
% (Hemigroup.levyExponentD_ne_top_of_integrableOn,
%  Hemigroup.SelfDecomposableExponent.integrableOn_of_ne_top)
\begin{lemma}[the integrability criterion for a profile]\label{lem:profile-integrability}
\uses{prop:fourier-toolbox}
\lean{SpatialLine.profile_integrability, SpatialLine.profile_integrability_mem}\leanok
Let $k : (0,\infty) \to [0,\infty)$ be measurable and put $\nu(dx) := k(x)\,x^{-1}dx$. Then
\[
  \int_{(0,\infty)} (1 \wedge x^2)\,\nu(dx) < \infty
  \qquad\Longleftrightarrow\qquad
  \int_0^1 x\,k(x)\,dx < \infty \ \ \text{and} \ \
  \int_1^\infty \frac{k(x)}{x}\,dx < \infty .
\]
Consequently, by \ref{prop:fourier-toolbox}(3), the right-hand pair of conditions is necessary
and sufficient for \ref{eq:sd-profile} to define a member of $\NDs$.
\statusT\quad\emph{(A change of variables and nothing else. It is stated because it is the
admissibility test every family in this article is put through, and because its two halves are
the symmetric forms of the causal criterion: near the origin the integrand is
$\asymp \omega^2x^2\,k(x)/x$, so $\int_0^1 x\,k$ replaces the causal $\int_0^1 k$, while at
infinity $1 - \cos \le 2$ gives the same condition as there. Machine-checked 2026-09-09. The
criterion itself is proved on Lean core alone, and takes $k$ measurable and nonnegative and
\emph{nothing else} --- in particular not monotone, because
\ref{lem:cin-rays} and \ref{prop:choquet-cone} apply it to increment densities
$k(u/t) - k(u/s)$, which are not. The second sentence is split in the Lean into the production
of the pair, which is again Lean core, and the membership of $\NDs$, which is where
\ref{prop:fourier-toolbox}(3) --- ledger A3, converse clause --- is spent; that split is what
makes the trust boundary of the two halves separately visible.)}
\end{lemma}

\begin{proof}\leanok
Both terms below are nonnegative, so
\[
  \int_{(0,\infty)}(1 \wedge x^2)\,\nu(dx)
  = \int_0^1 x^2\,k(x)\,\frac{dx}{x} + \int_1^\infty k(x)\,\frac{dx}{x}
  = \int_0^1 x\,k(x)\,dx + \int_1^\infty \frac{k(x)}{x}\,dx,
\]
and the left side is finite exactly when both terms are. The second sentence is
\ref{prop:fourier-toolbox}(3) applied to the pair $(a,\nu)$.
\end{proof}

% CHANGED (article mirror 2026-09-14, R157): the remark is split in two, as the paper prints it,
% and its first half is reworded. "Rate per unit of scale" is the linear-accumulation reading the
% article exists to drop -- the catalogue of one law is an intensity, and how the catalogue depends
% on scale is chapter 7's result, not a rate; "diffusive jitter" is out (the word was the author's
% objection, and "a Gaussian displacement of variance 2a" says what it is); and the Poisson stream
% is explained in its own sentence rather than named in a subordinate clause.
\begin{remark}[the displacement dictionary]\label{rem:displacement-dictionary}
A symmetric kernel is the law of a \emph{random displacement}: $\Phis{s}{t}f$ reports the
signal at a randomly displaced position $x - X_{s,t}$ with $X_{s,t} \sim \mu_{s,t}$. Its
transform $\hat\mu_{s,t}(\omega) = \mathbb{E}\cos(\omega X_{s,t})$ is the contrast that a
sinusoidal probe of wavenumber $\omega$ retains after the random displacement, the stage's
modulation transfer function (\ref{rem:modulation-transfer}). The exponent
$\psi = -\log\hat\mu$ adds along the cascade. In \ref{eq:levy-khintchine} the Gaussian
coefficient $a$ is the \emph{diffusive part}, a Gaussian displacement of variance $2a$, the one
term the time-causal axis deletes and the only one Gaussian scale space uses. The measure $\nu$
is a \emph{catalogue of jump displacements}. Displacements of size about $x$ arrive as a
Poisson stream, that is at independent random moments with a constant rate, and $\nu(dx)$ is
that rate; the streams of different sizes are independent, and each event of size $x$ costs the
probe the phase deficit $1 - \cos\omega x$. So a symmetric infinitely divisible displacement is
a Gaussian part plus an independent sum of jumps of all sizes, and the representation says that
every such displacement decomposes in this way.
\end{remark}

\begin{remark}[the profile in log-displacement]\label{rem:displacement-profile-log}
In the log-displacement coordinate $\theta = \log x$ the density shape of \ref{eq:sd-profile}
is the profile $\tilde k(\theta) = k(e^\theta)$ per unit log-displacement, dilation acts by
translation in $\theta$, and self-decomposability is the profile being nonincreasing: zooming
out may only add displacement intensity. The Gaussian is the profile-free member; the symmetric
stable exponents $|\omega|^\alpha$ are the exactly scale-free profiles
$\tilde k = c\,e^{-\alpha\theta}$. The extreme rays of the cone are the Gaussian ray $\omega^2$
and the unit steps, whose exponents are the $\Cin$ functions of \ref{lem:cin-rays}.
\end{remark}
